\documentclass[11pt]{article}
\usepackage[letterpaper,margin=0.86in]{geometry}
\usepackage[T1]{fontenc}
\usepackage{lmodern,amsmath,amssymb,amsthm,booktabs,microtype}
\usepackage{hyperref,fancyhdr}
\hypersetup{colorlinks=true,urlcolor=blue,linkcolor=blue,
pdftitle={Albertson--Collins: proof and verification},pdfauthor={Prepared for Alex Chan}}
\setlength{\parindent}{0pt}
\setlength{\parskip}{5pt}
\setlength{\headheight}{14pt}
\pagestyle{fancy}\fancyhf{}
\fancyhead[L]{Albertson--Collins: proof and verification}
\fancyhead[R]{28 October 2026}
\fancyfoot[C]{\thepage}
\newcommand{\Z}{\mathbb Z}
\newcommand{\F}{\mathbb F}
\newcommand{\Aut}{\operatorname{Aut}}
\newcommand{\Sym}{\operatorname{Sym}}
\newcommand{\AGL}{\operatorname{AGL}}
\begin{document}
{\LARGE\bfseries The canonical-coloring theorem}\\[6pt]
{\large Mathematical assessment and working guide for Alex Chan}\\
Harvard University and NBER \quad \texttt{achan@hbs.edu}

\section*{1. Assessment and exact scope}
\textbf{Assessment: I find no mathematical gap in the uploaded all-orders
argument.} With the classical Burnside--Schur theorem as an external
mathematical input, it proves the prescribed-coloring conjecture for
\emph{every} cyclic tournament. This is a mathematical review, not a
Lean kernel certificate or external human peer review. A bounded
literature search did not locate a later complete resolution; priority
is not established.

Write $n=2q+1\geq3$, and let $S$ contain exactly one of $d,-d$ for each
nonzero residue modulo $n$. The tournament $T(n;S)$ has an arc $i\to j$
when $j-i\in S$. The required coloring is
\[
 \lambda(i)=1\ (0\leq i\leq q),\qquad
 \lambda(i)=2\ (q<i\leq2q).
\]
It is \emph{distinguishing} when its only color-preserving
automorphism is the identity. It need not be a proper vertex coloring.

This is exactly Conjecture 2 in Meslem--Sopena [1]. The original
Albertson--Collins note [2, Conjecture 2.4] uses the smaller interval
$L=\{0,\ldots,q-1\}$. Its complement is
$\{q,\ldots,2q\}=I_n+q$, where $I_n=\{0,\ldots,q\}$.
A set and its complement have the same stabilizer, and translating
a set conjugates its stabilizer. Thus these formulations are equivalent.
The existence of \emph{some} distinguishing two-coloring was already
known from Gluck's theorem; the prescribed interval is the result at issue.

\textbf{Stronger theorem.} If $n$ is a positive odd integer and
$G\leq\Sym(\Z_n)$ has odd order and contains $\tau:x\mapsto x+1$,
then the only $g\in G$ satisfying $g(I_n)=I_n$ is the identity.

A \emph{regular subset} means a subset with trivial setwise stabilizer.
The theorem covers all connector sets without using their details.
There is an infinite family with $2^q$ connector sets at order $2q+1$,
so no finite collection of graphs alone can prove the statement.

\textbf{External input.} Burnside--Schur [3, opening theorem, p.~1]
says that a primitive action with a cyclic regular subgroup is either
doubly transitive or has prime degree and is conjugate to an affine
action over the prime field. The action-level conclusion is essential:
an abstract group isomorphism would not suffice.

\textbf{What changed in the new paper.} The supplied induction and
coordinate-normalization argument are retained. An elementary
interval-intersection lemma replaces the finite-field second-moment
calculation. Both versions have been checked.

\newpage
\section*{2. Full proof: the block-system step}
We prove the stronger theorem by strong induction on odd $n$. For
$n=1$, there is only the identity permutation. Assume the claim is
known for every smaller positive odd degree.

First record two consequences of the induction hypothesis. A set
and its complement have identical stabilizers. Also, if a group
contains all translations, then regularity of $I$ implies regularity
of $I+t$: a permutation $h$ preserving $I+t$ gives
$\tau^{-t}h\tau^t$ preserving $I$. These facts handle the smaller
half-intervals that appear inside blocks.

The group $G$ is transitive because it contains the full $n$-cycle
$\tau$. Suppose it has a nontrivial invariant block system.
Let $B$ be its block containing $0$. For $t\in B$, the point $t$
belongs to both $B$ and $B+t$. The block property gives $B+t=B$.
This implies closure under addition and also gives $-t\in B$,
since $0\in B+t$. Hence $B$ is an additive subgroup of $\Z_n$.

Every subgroup of $\Z_n$ consists of multiples of a divisor $a$
of $n$. Let $n=ab$. All blocks are translates of $B$, so they are
\[
 B_r=\{r+at:0\leq t<b\},\qquad r=0,\ldots,a-1.
\]
Nontriviality gives $1<a,b<n$, and both factors are odd. Write
$a=2u+1$ and $b=2v+1$. Since $(ab-1)/2=av+u$,
\[
 I_n\cap B_r=
 \begin{cases}
 \{r+at:0\leq t\leq v\},&r\leq u,\\
 \{r+at:0\leq t<v\},&r>u.
 \end{cases}
\]
Thus the intersection counts are $v+1$ exactly for block indices
$0,\ldots,u$, and $v$ for the others. These coordinates have
$0\leq r+at<ab$, so there is no wraparound in the calculation.

Take $g\in G$ with $g(I_n)=I_n$. Its induced permutation on blocks
preserves these counts, hence preserves $I_a$.
The full induced group is an image of $G$, so has odd order.
It contains $r\mapsto r+1$ on $\Z_a$.
Induction makes the induced permutation of $g$ the identity.
Consequently $g$ fixes each block \emph{setwise}.

Now restrict to $B_r$. The image on $B_r$ of its setwise stabilizer
in $G$ has odd order. It contains $\tau^a$, which is
$t\mapsto t+1$ in the coordinate $r+at\leftrightarrow t\in\Z_b$.
The preserved subset is either $I_b=\{0,\ldots,v\}$ or
$\{0,\ldots,v-1\}$. The latter has complement
$\{v,\ldots,2v\}=I_b+v$. Induction, complementing and translating
as appropriate, therefore makes the restriction of $g$ the identity.
This holds in every block, so $g$ is the identity.

\textbf{Example at $n=9$.} For $a=b=3$, the blocks are
$\{0,3,6\},\{1,4,7\},\{2,5,8\}$. Their intersections with
$I_9=\{0,1,2,3,4\}$ have sizes $2,2,1$. The quotient argument
fixes the blocks; the within-block argument fixes their points.

\newpage
\section*{3. Full proof: the primitive step and the tournament}
Suppose the action is primitive. It cannot be doubly transitive:
such a group has order divisible by the number $n(n-1)$ of ordered
distinct pairs, which is even for $n>1$ odd. Burnside--Schur
therefore gives prime degree $n=p$ and a bijection
$f:\Z_p\to\F_p$ conjugating $G$ into $\AGL_1(p)$.

\textbf{The original labels are retained.} The permutation
$f\tau f^{-1}$ has order $p$. An affine map $z\mapsto\alpha z+\beta$
with $\alpha\neq1$ has a fixed point; moving that point to $0$
conjugates it to multiplication by $\alpha$. Its order divides
$p-1$. Hence $f\tau f^{-1}$ is translation by some $c\neq0$, and
\[
 f(j+1)=f(j)+c,\qquad f(j)=f(0)+jc.
\]
Renormalizing by $z\mapsto(z-f(0))/c$ gives the original residue
coordinate. Every element of $G$ is therefore affine
\emph{in the coordinates defining $I_p$}.

\textbf{Affine interval lemma.} Let $p=2q+1$ and
$I=\{0,\ldots,q\}$. For a residue $d$ put
$C(d)=|I\cap(I+d)|$. For $1\leq d\leq q$,
$I\cap(I+d)=\{d,\ldots,q\}$; also $C(-d)=C(d)$ by translation.
Thus
\[
 C(0)=q+1,\qquad
 C(d)=q+1-\min\{d,p-d\}\quad(1\leq d\leq2q).
\]
The largest value at nonzero shifts is $q$, occurring exactly at
$d=\pm1$. If $h(x)=\alpha x+\beta$ preserves $I$, then it maps
$I+d$ to $I+\alpha d$, giving $C(\alpha d)=C(d)$.
Taking $d=1$ forces $\alpha=\pm1$.

If $\alpha=1$, a nonzero $\beta$ would imply
$C(\beta)=|I|=q+1$, impossible. Thus $h$ is the identity.
The reflection $r(x)=q-x$ preserves $I$.
If $\alpha=-1$, then $r\circ h$ is a translation preserving $I$,
so it is the identity and $h=r$. This is a nonidentity involution,
excluded by odd group order. The primitive case, and the induction,
are complete.

The affine lemma works even at $p=3$. It also holds over every odd
composite modulus when $\alpha$ is a unit, although the primitive
case only needs a prime modulus.

\textbf{Tournament corollary.} A finite tournament has odd-order
automorphism group. Otherwise Cauchy's theorem gives a nonidentity
involution. Such a permutation swaps some distinct $u,v$, reversing
the unique arc joining them, and therefore cannot be an automorphism.
A cyclic tournament has $\tau$ among its automorphisms because
translation preserves all differences. Apply the group theorem to
$\Aut(T)$. Preserving the two colors is precisely preserving
$I_{2q+1}$ setwise, so only the identity can do so.
This proves the conjecture for every $q$ and every connector set.

\newpage
\section*{4. The supplied proof and computational verification}
The uploaded second-moment argument is also valid. For $p\geq5$,
put $m=(p+1)/2$ and work in $\F_p$. The centroid of $I_p$ is
$c=(m-1)/2=-1/4$. The sum-of-squares identity gives
\[
 M=\sum_{j=0}^{m-1}(j-c)^2
   =\frac{m(m^2-1)}{12}
   =-\frac1{32}\neq0.
\]
All denominators are invertible for $p\geq5$. A preserving affine
map fixes $c$ and multiplies $M$ by $\alpha^2$, so $\alpha^2=1$.
The positive multiplier gives the identity; the negative multiplier
gives an involution. The supplied proof treats $p=3$ separately,
avoiding division by $3$ there.

\textbf{Actual graphs.} The Python module constructs any input
$T(n;S)$, validates $S$, and exports arcs, adjacency matrices,
directed DOT files, and SVG diagrams. The included examples are:
\begin{center}
\begin{tabular}{@{}lrcrr@{}}
\toprule
Example & $n$ & $S$ & $|\Aut(T)|$ & Colored order\\
\midrule
Directed triangle & 3 & $\{1\}$ & 3 & 1\\
Consecutive connectors & 5 & $\{1,2\}$ & 5 & 1\\
Paley tournament & 7 & $\{1,2,4\}$ & 21 & 1\\
Three triangles in a cycle & 9 & $\{1,3,4,7\}$ & 81 & 1\\
\bottomrule
\end{tabular}
\end{center}
The final example is the imprimitive composition $C_3[C_3]$.
Checking translations alone would miss many of its automorphisms.
The implemented search uses arbitrary adjacency-preserving
permutations, not an assumed affine family.

\textbf{Checks completed.}
\begin{itemize}
\item All $4{,}094$ connector sets of orders $3,5,\ldots,23$, plus
the singleton control, passed exact canonical-coloring checks.
\item An independent all-permutations oracle matched the
automorphism engine in $2{,}258$ cases.
\item All $143{,}346$ affine permutations over the 50 odd moduli
from 3 to 101 passed the affine interval lemma check.
\item The supplied sanity program was rerun: 167 odd primes below
1,000, 2,401 factor pairs, and all relevant odd-order color-preserving
permutations through degree 13 passed.
\end{itemize}
These finite checks support the formulation and code; they are
\emph{not premises} of the all-orders proof. The unrelated conjecture
searches in the Round41 archive are not used and have not been
re-certified here.

\textbf{Provenance.} The separately attached PDF and the PDF inside
the archive have identical text after whitespace normalization.
The project includes input hashes, an independent proof audit, the
new affine-lemma check, and machine-readable finite-check results.

\newpage
\section*{5. Formalization and publication}
\textbf{Lean status.} The supplied Lean files are a foundation,
not a complete formal proof. Lean and Lake are unavailable in this
execution environment, so no file has been compiled or kernel-checked.
Comments and the Lean README distinguish proposed proofs, conditional
interfaces, and remaining mathematical obligations.

The main remaining work is the universal group theorem: classify
translation-invariant blocks, construct quotient and within-block
actions, prove the interval counts, and carry out strong induction.
The primitive step needs the action-level Burnside--Schur theorem,
coordinate normalization, and the affine interval lemma.
Absence of \texttt{sorry} alone does not make a conditional theorem
unconditional. A future certificate should expose dependencies with
\texttt{\#print axioms}; treating Burnside--Schur as an unproved axiom
would produce a relative certificate.

\textbf{Paper status.} The accompanying short article contains the
full proof and the supplied author information. It is a submission
draft, with no claim of journal acceptance, Lean verification, or
established priority. The current literature search checked primary
papers and targeted later work, but did not exhaust MathSciNet,
zbMATH, or unpublished work. An independent human reading and a
citation-forward novelty check remain useful before submission.
Journal-specific preparation notes are included separately.

\textbf{Primary references checked.}
\begin{enumerate}
\item K. Meslem and E. Sopena, \emph{On the distinguishing number of
cyclic tournaments: Towards the Albertson--Collins Conjecture},
Discrete Applied Mathematics \textbf{266} (2019), 219--236.
Conjecture 2; \href{https://arxiv.org/html/1608.04866v4}{arXiv:1608.04866v4}.
\item M. O. Albertson and K. L. Collins, \emph{A note on breaking the
symmetries of tournaments}, Congressus Numerantium \textbf{136}
(1999), 129--131. Conjecture 2.4;
\href{https://kcollins.web.wesleyan.edu/publications/symnote-saks.pdf}{coauthor's primary copy}.
\item S. Evdokimov and I. Ponomarenko, \emph{A new look at the
Burnside--Schur theorem}, Bulletin of the London Mathematical Society
\textbf{37} (2005), no. 4, 535--546. Opening theorem, p. 1;
\href{https://arxiv.org/pdf/math/0310373}{arXiv:math/0310373}.
\item D. Gluck, \emph{Trivial set-stabilizers in finite permutation
groups}, Canadian Journal of Mathematics \textbf{35} (1983), no. 1,
59--67. Corollary 1;
\href{https://doi.org/10.4153/CJM-1983-005-2}{doi:10.4153/CJM-1983-005-2}.
\end{enumerate}

\textbf{Continue from this package.} Compile the elementary Lean
layer in a pinned Lean/mathlib environment, then replace the
conditional group-theorem interface with the induction proof.
The project README gives the source files and reproduction commands.
\end{document}
